\documentclass[11pt]{article}
\usepackage[margin=1.1in]{geometry}
\usepackage{amsmath,amssymb,amsthm}
\usepackage{booktabs}
\usepackage{lmodern}
\usepackage{microtype}
\usepackage[colorlinks=true,linkcolor=blue,citecolor=blue,urlcolor=blue]{hyperref}

\newtheorem{conjecture}{Conjecture}
\newtheorem{theorem}{Theorem}
\newtheorem{lemma}{Lemma}
\newtheorem{remark}{Remark}
\theoremstyle{definition}
\newtheorem{question}{Question}

\newcommand{\Li}{\operatorname{Li}}
\newcommand{\e}{\mathrm{e}}
\DeclareMathOperator{\ord}{ord}
\newcommand{\SG}{\mathfrak S}
\newcommand{\CT}{C_3}
\newcommand{\M}{\mathcal M}

\title{Twenty-five conjectures from a local--global random model}
\author{}
\date{July 27, 2026}


\begin{document}

\title{Twenty-five conjectures from a local--global random model:\\ the statements}
\author{}\date{}
\maketitle
\section{Notation}

$p,q$ denote primes; $\gamma$ is Euler's constant,
$\e^{\gamma}=1.781072\ldots$; $\varphi$ is Euler's totient;
$\phi=(1+\sqrt5)/2$; $F_n$ the Fibonacci numbers; $p\#$ the primorial.
For irreducible $f_1,\dots,f_k\in\mathbb Z[x]$ with positive leading
coefficients, the Bateman--Horn singular series is
\[
  C(f_1,\dots,f_k)\;=\;\prod_{p}\frac{1-\omega(p)/p}{(1-1/p)^{k}},
  \qquad
  \omega(p)=\#\{n \bmod p:\ p\mid f_1(n)\cdots f_k(n)\},
\]
and the corresponding main term is
\[
  I(N)\;=\;\int_{2}^{N}\frac{dt}{\prod_{i}\log f_i(t)},
\]
which absorbs the usual $1/\prod_i \deg f_i$.  The Bateman--Horn
conjecture \cite{BH} asserts
$\#\{n\le N: \text{all } f_i(n) \text{ prime}\}\sim C\cdot I(N)$
whenever $C\neq0$; each Bateman--Horn conjecture below is an instance
with its constant evaluated.  For nonlinear systems the product
converges only conditionally; throughout, it is taken in the canonical
ordering of increasing primes, in which convergence follows from the
equidistribution of $\omega(p)$ about its Chebotarev mean.  Our
numerical truncation wobble is an error bar for this canonically
ordered product.  We write
$\Li_k(x)=\int_2^x(\log t)^{-k}\,dt$.  The twin singular series is
$\SG(d)=2C_2\prod_{p\mid d,\,p>2}\frac{p-1}{p-2}$ for even $d$, with
$2C_2=1.3203236\ldots$; the Goldbach series $\SG(n)$ is the same product
over odd $p\mid n$.

Eight of the conjectures below are built over the Bateman--Horn
move---choose an admissible system, compute its singular series,
state the count---but none is a bare instance of it:
\ref{c1} and \ref{c2} are family laws with derived constant
distributions, \ref{c4} and \ref{c5} carry structural classifications,
\ref{c6} is a searched chain, and \ref{c3}, \ref{c7}, \ref{c8} are
contamination-calculus predictions (two fresh races and the negative
control) whose mechanism is the contamination calculus of
Conjecture~\ref{c21}(v).  Each carries
its classical calibration instance as an attributed remark.
The heuristic is uniform, so we give it once.  Model the primality of the
values $f_i(n)$ as independent events
of probability $1/\log f_i(n)$, corrected by the factor
$(1-\omega(p)/p)/(1-1/p)^k$ at each prime $p$ to account for the joint
local behaviour; multiplying the corrections and integrating gives
$C\cdot I(N)$.

For locally integrable $F$ we write
$\M_x(F)=\frac1{\log x}\int_2^x F(t)\,\frac{dt}t$ for the logarithmic
mean; all race conjectures below state their drift and null clauses
through $\M_x$, so that ``no persistent leader'' and ``systematic
surplus'' are claims about a defined functional rather than about
pointwise behaviour.  Drift clauses are stated at the \emph{drift
scale}---for a pair race, $\M_x(D(t)\log^2t/\sqrt t)$ converging to
an explicit constant---never as
$\M_x((D-T)/\sqrt\pi)\to0$, which is vacuous for the drift coefficient
(both $T/\sqrt\pi$ and its multiples have logarithmic mean $\to0$).
Convergence of $\M_x$ at
drift scale is itself part of each such conjecture: a pure
random-walk fluctuation model would make these means divergent in
variance, so their convergence encodes the Rubinstein--Sarnak-type
almost-periodicity that the zero-oscillation hypothesis names.

Admissibility ($\omega(p)<p$ for all $p$) is checked for every constant
computed.  The natural-looking pair $\{n^2+n+1,\ n^2+n+3\}$ has
$\omega(3)=3$ and is
rejected; the $k=4$ branch of Conjecture~\ref{c5} and the naive
iterated chain behind Conjecture~\ref{c6} are inadmissible for the
same reason, and both exclusions are part of the statements.


\section{Principal conjectures}

\begin{lemma}[Orientation elimination and class assignment]\label{l:c21}
In $\psi_2(x)=\sum_{n\le x}\Lambda(n)\Lambda(n+2)$, the total
contribution of terms in which $n$ or $n+2$ is a proper prime power
is, apart from $O(x^{1/3}\log^2 x)$ from cubes and higher powers,
carried entirely by the patterns $(q^2-2,\,q^2)$ with $q$ prime and
$q^2-2$ prime: the mirror pattern $(q^2,\,q^2+2)$ is annihilated by
$3\mid q^2+2$ for every prime $q>3$.  Moreover, for every prime
$q\neq5$, $q^2-2\equiv2$ or $4\pmod 5$ (never $1$), and for every odd
$q$, $q^2-2\equiv7\pmod8$.  (The single exceptional term $q=5$,
$(23,25)$, lands outside the twin-start classes altogether and is
absorbed in the bounded error.)
\end{lemma}

\setcounter{conjecture}{0}%
\begin{conjecture}[Twin races modulo 5 and 8, and the prime-power
contamination calculus; apparently new]\label{c21}
Twin starts $p>5$ lie in classes $p\equiv1,2,4\pmod5$.  Let
$D_1(x)=\pi_t(x;5,1)-\tfrac12\bigl(\pi_t(x;5,2)+\pi_t(x;5,4)\bigr)$
and
\[
T(x)\;=\;\frac1{2\log^2x}
\sum_{\substack{q\le\sqrt x\\ q^2-2\ \mathrm{prime}}}
\log q\,\log(q^2-2),
\]
whose scale is governed by the constant $C_7$ of
Conjecture~\ref{c7}.  Then:
\emph{(i)} [drift law, at drift-scale normalization] class $1$
carries the systematic surplus $T$.  A normalization at the noise scale,
$\M_x((D_1-T)/\sqrt{\pi_t})\to0$, would be \emph{vacuous} for the
coefficient: since $T/\sqrt{\pi_t}\asymp1/\log t$, one has
$\M_x(T/\sqrt{\pi_t})\asymp\log\log x/\log x\to0$, so the condition
holds with $T$ replaced by $0$, $2T$, or any fixed multiple, and
does not express ``exactly the removed mass.''  The statement
therefore normalizes at the \emph{drift} scale, where the claimed
limit is a nonzero constant:
\[
\M_x\!\left(\frac{D_1(t)\,\log^2t}{\sqrt t}\right)
\;\longrightarrow\;
c_T\;:=\;\lim_{x\to\infty}\frac{T(x)\log^2x}{\sqrt x}
\;=\;\tfrac12\,C(x,\,x^2-2),
\]
the half of the Bateman--Horn constant of the pair $(q,\,q^2-2)$
(the weighted census sum is $\sim C\sqrt x$, so the limit exists and
is explicit).  Two layers are being conjectured at once: \emph{that} the
logarithmic mean converges at this
normalization---under a pure random-walk model it would not (its
variance grows like $\log x$), so convergence is an assertion of
Rubinstein--Sarnak-type almost-periodic structure for the twin race,
and is the precise content of the zero-oscillation
hypothesis---and that its limit is $c_T$ rather than $0$ or $2c_T$,
which is the calculus.  Equivalently:
$\M_x(D_1)/\M_x(T)\to1$;
\emph{(ii)} [zero deterministic drift in the symmetric differences]
the symmetric classes carry no deterministic term \emph{at the same
normalization}:
$\M_x\bigl((\pi_t(\cdot;5,2)-\pi_t(\cdot;5,4))\log^2t/\sqrt
t\bigr)\to0$,
and likewise for each of the three pairwise differences among classes
$\{1,3,5\}$ mod $8$ in clause (iv)---now a statement distinct from (i),
since at drift scale ``zero'' and ``$c_T$'' are
different limits;
\emph{(iii)} [sign density, conditional] conditional on the Gaussian
fluctuation model for the normalized remainder in (i)--(ii), the
logarithmic density of $\{D_1>0\}$ exists and equals $\tfrac12$,
approached from above with a finite-$x$ excess of order
$\log\log x/\log x$: the pointwise drift-to-noise ratio is
$T/\sqrt{\pi_t}\asymp1/\log t$, and its \emph{logarithmic average}
$\M_x(1/\log t)\sim\log\log x/\log x$ carries an extra $\log\log$
(the scale is not $1/\log x$).  In contrast to Chebyshev's races, the twin
race has no
persistent leader;
\emph{(iv)} [mod-8 companion] since $q^2-2\equiv7\pmod8$ always
(Lemma~\ref{l:c21}), the \emph{entire} deterministic term falls on
the single class $7\pmod 8$: with
$D_7(x)=\tfrac13\sum_{a\in\{1,3,5\}}\pi_t(x;8,a)-\pi_t(x;8,7)$,
\[
\M_x\!\left(\frac{D_7(t)\,\log^2t}{\sqrt t}\right)
\longrightarrow2c_T
\]
(twice the mod-5 constant, undiluted---at this normalization the factor $2$
is a falsifiable prediction), while classes $1,3,5$ are mutually
symmetric.  The mod-5 and mod-8 races are two projections of one
mechanism, so their joint behaviour is a family test in the sense of
the shape-of-agreement criterion;
\emph{(v)} [contamination calculus for balanced races] for a pattern
$(n,n+d)$ and modulus $m$, call the race \emph{balanced} if the
prime-prime singular series is identical across the competing residue
classes (as CRT gives for the classes compared in (i)--(iv)) and no
mechanism other than prime powers contributes at the $\sqrt x$ scale
under the zero-oscillation hypothesis.  For balanced races the
deterministic drift vector is computed by one operator.  Define the
\emph{contamination operator}
\[
\mathcal C_{d,m}(a;x)\;=\;
\sum_{\text{orientations }o\in\{(q^2-d,q^2),\,(q^2,q^2+d)\}}
\ \sum_{\substack{q\le\sqrt x,\ o\ \text{prime-compatible}\\
\mathrm{start}(o)\equiv a\ (m)}}
\Lambda(q^2)\,\Lambda(\text{prime member of }o),
\]
where prime-compatible means the non-square member is prime and no
fixed prime annihilates the orientation.  The provable layer is the
orientation census and class assignment (which $o$ survive, and
where their starts land---Lemma~\ref{l:c21} and its analogues); the
conjectural layer is the transfer: the unweighted class counts are
depressed by $\mathcal C_{d,m}(a;x)/\log^2x$ relative to the clean
classes, in the drift-scale limiting-log-mean sense of (i)---for
each class $a$, the deficit obeys
$\M_x\bigl(\mathrm{deficit}_a(t)\log^2t/\sqrt t\bigr)\to
c(a):=\lim\mathcal C_{d,m}(a;x)/\sqrt x$, an explicit Bateman--Horn
constant per class, so every coefficient in the vector is
identified---i.e.\ ``exactly the removed
mass,'' which requires the weighted identity, partial summation
uniform in classes, higher-power bounds (exponents $r\ge3$
contribute $O(x^{1/3+\varepsilon})$ against the operator's
$\asymp\sqrt x$ and are negligible---the one easy layer of the
transfer), and the zero-oscillation
hypothesis, and is the calculus's single analytic conjecture rather than a
rule.  The drift is
$\asymp\sqrt x/\log^2x$ against a count fluctuation
$\asymp\sqrt x/\log x$, i.e.\ smaller by $1/\log x$ in
standard-deviation units at every height---so the calculus predicts
occupation biases under long logarithmic aggregation, never a fixed
nonzero standardized mean, and every verification in this family is
designed around that fact.  (Races with unequal primary
constants, or with interfering orientations beyond those enumerated,
are outside the calculus until those terms are separated.)  In
particular,
for the cousin pattern $d=4$ the calculus makes two predictions with
no free choices: the orientation $(q^2-4,q^2)$ is dead
($q^2-4=(q-2)(q+2)$ is composite), the orientation $(q^2,q^2+4)$
survives (no prime annihilates it), $q^2+4$ prime forces
$q\equiv\pm2\pmod5$ for $q\neq5$ (at $q=5$ the value $29$ is prime,
but the start $25\equiv0\pmod5$ lies outside the cousin start classes
$\{2,3,4\}$), so the contaminated cousin-start class is
$4\pmod 5$ (of $\{2,3,4\}$) and---since $q^2\equiv1\pmod8$---the
class $1\pmod 8$ (of $\{1,3,5,7\}$): both deficits converge, in the
drift-scale limiting-log-mean sense of (i)--(ii), to
$c^{\mathrm c}=\lim_{x\to\infty}T^{\mathrm c}(x)\log^2x/\sqrt x$,
the Bateman--Horn constant of the pair $(q,\,q^2+4)$, where
$T^{\mathrm c}(x)=\frac1{\log^2x}\sum_{q\le\sqrt x,\ q\neq5,\ q^2+4\
\mathrm{prime}}\log q\,\log(q^2+4)$.
\end{conjecture}

\setcounter{conjecture}{1}%
\begin{conjecture}[The pair-level Montgomery--Soundararajan reduction;
apparently new]\label{c12}
For the window field of Conjecture~\ref{c16}(ii), the diagonal
variance of the twin-pair count reduces to a \emph{pinned}
singular-series average.  With
$R(h)=C_4(0,2,h,h+2)/\SG(2)^2$ for $h\ge1$ (zero when the offset set
is degenerate or inadmissible---all odd $h$, and $h=2$ since
$\{0,2,4\}$ dies at $3$; consecutive twin pairs cannot overlap beyond
$(3,5,7)$, so no overlap term exists) and
\[
G(H)\;=\;\sum_{1\le h\le H}\Bigl(1-\frac hH\Bigr)\bigl(R(h)-1\bigr),
\]
the derived identity is
\[
\frac{\operatorname{Var}_t\,\pi_2(t;H)}{\mathbb E_t\,\pi_2(t;H)}
\;=\;1+\frac{\SG(2)\,\bigl(2G(H)-1\bigr)}{\log^2x}
+o\!\Bigl(\frac1{\log^2 x}\Bigr):
\]
with $p=\SG(2)/\log^2x$ the per-site intensity, the variance is
$Hp(1-p)+2p^2HG(H)$, so $\operatorname{Var}/\mathbb E=1-p+2pG(H)$
---the Bernoulli diagonal term $-p=-\SG(2)/\log^2x$ is of exactly
the order of the claimed error and cannot be dropped (it is the same
diagonal that, in the single-prime case, merges into
Montgomery--Soundararajan's constant $\gamma+\log2\pi-1$).  With
that term retained, the entire deficit left open at
Conjecture~\ref{c16}(ii) is
the single computable function $G$ plus the explicit Bernoulli
correction---the pair analogue of
Montgomery--Soundararajan's reduction of the prime-count variance to
$\sum_{h\le H}\SG(h)(H-h)$ \cite{MS}.  Clause (i) is a
\emph{conditional proposition}, not an independent conjecture: given
quantitative Hardy--Littlewood for $4$-tuples, the reduction
follows.  The conjectural clause is (ii): $-G(H)/\log H$ diverges
---the pinned average grows strictly faster than the single-prime
secondary term $\tfrac12(\log H+\gamma+\log2\pi-1)$, so pair counts
in windows are more sub-Poisson than prime counts---with the caveat that
data to $H=3000$ cannot distinguish divergence from a pure
$\log$ with a large constant
(nor $c\log H\log\log H$ from $(\log H)^\alpha$): what the
computation \emph{establishes} is only that the local slope exceeds
the single-prime value sevenfold across the computable range; the
divergence and its exact form are registered, not evidenced beyond
that.
\end{conjecture}

\setcounter{conjecture}{2}%
\begin{conjecture}[Least primes in progressions; part (i) previously
stated \cite{HLS,Wagstaff,Fiori,Leung}, part (ii) apparently
new]\label{c22}
Let $p(a,q)$ be the least prime $\equiv a\pmod q$ and
$U(a,q)=\Li(p(a,q))/\varphi(q)$.  Then:
\emph{(i)} $U\Rightarrow\mathrm{Exp}(1)$ marginally as $q\to\infty$;
and, under the additional hypothesis that the class processes are
asymptotically jointly independent in the extreme-value sense (which
marginal convergence alone does not supply),
$\max_a U-H_{\varphi(q)}$ converges to a Gumbel law
($H_n=\sum_{k\le n}1/k$), the order-statistics form of Wagstaff's
$\max_a p(a,q)\sim\varphi(q)\log^2q$;
\emph{(ii)} [canonical deficit, dyadic-average form] define, with no
reference to any control model,
$\Theta(q)=\bigl(1-\mathbb E_a[U(a,q)]\bigr)\log q$.  The conjecture
is stated for dyadic averages over prime moduli---pointwise
convergence along every prime $q$ may be too strong, since
arithmetic fluctuations in $q$ could persist:
\[
\frac1{\#\{q\in[Q,2Q]\ \text{prime}\}}
\sum_{\substack{Q\le q\le2Q\\ q\ \mathrm{prime}}}\Theta(q)
\;\longrightarrow\;\Theta\;>\;0
\qquad(Q\to\infty).
\]
(The formal claim is model-free: existence and strict
positivity of the dyadic-average limit.  The stronger reading---that
$\Theta$ exceeds the discreteness level any parity-aware control
produces---is the diagnostic layer below, not part of the display,
since a control constant has no place in a model-free statement.)
The limiting \emph{distribution} of $\Theta(q)$ over the dyadic
block is left open, and in particular we do \emph{not} assert
$\operatorname{sd}(\Theta(q))\to0$: fluctuations driven by the
arithmetic of $q$ and of $q-1$ (through the singular-series pair
terms $\SG(kq)$ and the class structure mod $q$) may well persist at
bounded size, so $\Theta(q)$ need not concentrate around its
dyadic mean.
\end{conjecture}

\setcounter{conjecture}{3}%
\begin{conjecture}[The Goldbach lane race; apparently new]\label{c25}
For $n\equiv2\pmod4$ let $R_1(n)$, $R_3(n)$ be the ordered Goldbach
representation counts in the lanes $p\equiv q\equiv1$ and
$p\equiv q\equiv3\pmod4$ (Conjecture~\ref{c15}), and
$D(n)=R_3(n)-R_1(n)$.  Prime-square terms $n=q^2+p$ ($q$ odd prime)
land \emph{exclusively} in the $(1,1)$ lane, since
$q^2\equiv1\pmod4$ and then $n-q^2\equiv1\pmod4$ automatically.
Consequently the $(3,3)$ lane leads on average: with
\[
D_{\mathrm{sys}}(n)\;=\;\frac{2}{\bar\ell(n)}
\sum_{\substack{q \text{ odd prime},\ q\le\sqrt n\\ n-q^2 \text{ prime}}}
\log q\,\log(n-q^2),
\qquad
\bar\ell(n)\;=\;\frac{n-6}{\displaystyle\int_3^{n-3}
\frac{dt}{\log t\,\log(n-t)}},
\]
(the analytic mean of $\log p\log(n-p)$ under the lane profile; the
empirical lane mean is its estimator).  Write $\mathbb E_x$ for the
average over $n$ drawn from $\{n\equiv2\ (4),\ n\le x\}$ with weight
proportional to $1/n$ (the logarithmic sampling measure, now
explicit).  Then:
\emph{(i)} [drift law]
$\mathbb E_x[D-D_{\mathrm{sys}}]=o(\mathbb E_x[D_{\mathrm{sys}}])$ as
$x\to\infty$: the $(3,3)$ lane leads on average by exactly the square
contamination;
\emph{(ii)} [weighted mean and internal null lane]
$\mathbb E_x[D_{\mathrm{sys}}]$ is given by the Hardy--Littlewood
(Conjecture H) local model for $n-q^2$, whose local factors vary
strongly with $n$---most drastically at $p=3$: for $n\equiv1\pmod3$, $3\mid
n-q^2$ for \emph{every} prime
$q>3$, so the contamination collapses to exactly two boundary
families---the $q=3$ term (when $n-9$ is prime) and---a second family, not a
single term---the sparse exceptional family $n-q^2=3$ (when $n-3$ is
itself a
prime square, so that the divisible-by-3 value \emph{is} 3)---both
negligible in logarithmic mean, so the race carries no drift at the
systematic scale on the subprogression $n\equiv10\pmod{12}$: with the
positive comparison scale
$A_G(x):=\mathbb E_x^{(\not\equiv1\,(3))}[D_{\mathrm{sys}}]$, the
logarithmic-ensemble mean of the systematic term over the
contaminated lane $n\not\equiv1\pmod3$ (a comparison against this
subprogression's own systematic term, which vanishes, would be
empty), the claim is $\mathbb E_x[D]=o\bigl(A_G(x)\bigr)$ on
$n\equiv10\pmod{12}$, while the drift
concentrates on $n\not\equiv1\pmod3$: an internal null lane,
predicted by the same mechanism that predicts the lead;
\emph{(iii)} [sign density] the per-$n$ drift-to-noise ratio
$\kappa(n)=D_{\mathrm{sys}}(n)\big/\sqrt{\SG(n)J(n)}$
(where $J(n)=\int_3^{n-3}dt/(\log t\log(n-t))$, so that
$\SG(n)J(n)$ is the total ordered representation count) satisfies
$\kappa(n)\asymp c(n)/\log n\to0$; under the Gaussian fluctuation
model the logarithmic density of $\{D>0\}$ is therefore $\tfrac12$ in
the limit, with the finite-$x$ sign fraction predicted to be
$\mathbb E_x[\Phi(\kappa)]$ ($\Phi$ the standard normal distribution
function)---decaying, but far from $\tfrac12$ at any accessible
height.
\end{conjecture}

\setcounter{conjecture}{4}%
\begin{conjecture}[The least Goldbach summand: exponential law and
ordering deficit; the time-changed law and deficit constant
apparently unstated]\label{c15}
Let $s(n)$ be the least prime $p\nmid n$ with $n-p$ prime (the
exclusion $p\nmid n$ removes a trivial obstruction: $p\mid n$ forces $p\mid
n-p$, so $n-p$ is
composite except in the diagonal case $n=2p$, which would otherwise
distort the hazard at one atypical point), and
time-change by the expected-arrivals clock
$U(n)=\SG(n)\sum_{p\le s(n),\,p\nmid n}1/\log(n-p)$.  Then:
The ensemble is fixed first, since $U(n)$ is deterministic at each
$n$ and an expectation therefore needs a declared measure: for a scale
$X$,
\[
\mathbb E_XU=
\frac{\sum_{X<n\le2X,\ 2\mid n}U(n)/n}
{\sum_{X<n\le2X,\ 2\mid n}1/n},
\qquad
\Theta_G(X)=\bigl(1-\mathbb E_XU\bigr)\log X.
\]
\emph{(i)} [limit law] under $\mathbb E_X$-sampling,
$U\Rightarrow\mathrm{Exp}(1)$ as $X\to\infty$;
\emph{(ii)} [canonical ordering deficit]
$\Theta_G(X)\to\Theta_G>0$---the Goldbach sibling of the
least-prime deficit of Conjecture~\ref{c22}(ii), produced by the
same occupancy mechanism: the ordered primes fill the available
representations faster than an exchangeable model, so the first
arrival comes \emph{early} and $\mathbb E_X[U]<1$.
\end{conjecture}

\setcounter{conjecture}{5}%
\begin{conjecture}[{$n^2+2^n$}; sequence is OEIS A064539]\label{c11}
Every prime of the form $n^2+2^n$ with $n>1$ has $n\equiv3\pmod6$
(elementary), and:
\emph{(i)} infinitely many such primes exist, with counting function
$\asymp\log N$;
\emph{(ii)} [candidate law, entanglement-aware form] for a finite set
$S$ of primes $p\ge5$ (the primes $2$ and $3$ are exact in the lane
reduction, $\delta_2=\delta_3=0$, and $\ord_2 2$ is undefined) let
$D_S$ be the \emph{exact} joint density,
within the lane $n\equiv3\ (6)$, of $n$ with $p\nmid n^2+2^n$ for all
$p\in S$, computed over the full joint period
$\mathrm{lcm}_{p\in S}\,\mathrm{lcm}(6,\,p\cdot\ord_p2)$, and set
$\kappa_S=D_S/\prod_{p\in S}(1-1/p)$.  The conjecture is stated
along the \emph{canonical exhaustion} $S_z=\{p\ \text{prime}:5\le
p\le z\}$, and its three layers are logically separate and are
asserted separately:
\emph{(ii-a)} the sequence $\kappa_{S_z}$ converges as
$z\to\infty$ (a general net over arbitrary finite $S$ carries an
order-of-exhaustion ambiguity we do not need), with limit $\kappa$;
\emph{(ii-b)} whether the limit \emph{factors}---i.e.\ whether the
exact finite-level factorization observed below persists, so that
entanglement is asymptotically absent---is an \emph{open question},
expressly not counted among this paper's conjectural assertions, on which
(ii-a) takes no position; and
\emph{(ii-c)} the count is
$\sim\sum_{n\le N,\,n\equiv3\,(6)}\kappa/\log(n^2+2^n)$---a genuine
further step, since a convergent local density does not
automatically govern the global count.  The working
value $\kappa^{*}=4.2734\ldots$ is a \emph{hybrid}, not a
computation of any single $\kappa_{S_z}$: the CRT-exact core
$S_{19}$ times independent per-prime factors for $19<p\le300$---so
$\kappa^{*}=\kappa_{S_{300}}$ exactly if and only if the observed
exact factorization persists to $300$.
\end{conjecture}

\setcounter{conjecture}{6}%
\begin{conjecture}[Fermat quotients: the multibase subtorus law and
the Wieferich ledger; single-base heuristic previously stated
\cite{CDP}, the vertical multibase law apparently new]\label{c23}
With $q_p(a)=(a^{p-1}-1)/p\bmod p$ (defined for primes $p\nmid a$;
all clauses below range over such $p$ only) and $q_p=q_p(2)$, five
clauses:
\emph{(0)} [structure, classical] $q_p(ab)\equiv q_p(a)+q_p(b)
\pmod p$ (the Eisenstein--Lerch homomorphism; verified exactly on
every tested prime), so for multiplicatively \emph{dependent} bases
the vector $(q_p(a_1),\dots,q_p(a_r))/p$ is confined to the rational
subtorus cut out by the relations;
\emph{(iv)} [multibase equidistribution, the claimed law] for
multiplicatively \emph{independent} bases the vector equidistributes
on the full torus as $p$ varies---the \emph{vertical} joint law,
complementing the fixed-$p$, varying-base statistics of
Ostafe--Shparlinski and Cobeli--Zaharescu---with correlations
vanishing and the same LIL calibration as (i);
\emph{(v)} [simultaneous Wieferich] the expected count of $p\le x$
with $q_p(2)=q_p(3)=0$ has convergent sum $\sum1/p^2$: at most
finitely many simultaneous Wieferich primes exist (the single-base
folklore of this accounting is Conrad's), and the empirical list is
empty to $10^7$.  Three further clauses on the base-2 quotient
alone:
\emph{(i)} [global equidistribution] $q_p/p$ equidistributes on
$[0,1)$ with discrepancy at exactly the calibrated random-model
scale: in its sharp form,
\[
\limsup_{x\to\infty}
\frac{\sqrt{\pi(x)}\,D_{\mathrm{KS}}(x)}
{\sqrt{\log\log\pi(x)}}
\;=\;\frac1{\sqrt2},
\]
the Chung--Smirnov law-of-the-iterated-logarithm constant for an
i.i.d.\ uniform sequence---demanding both no more (a bound
$O(\pi(x)^{-1/2})$ would be \emph{stronger than the random model
supports}, the same over-demand that killed the Mertens conjecture)
and no less than the model earns: the upper bound alone would be consistent
with
hidden rigidity, and the matching lower bound asserts the quotients
fluctuate like genuine noise;
\emph{(ii)} [shrinking targets] for every fixed $K\ge1$, and uniformly
for $K=K(x)\le\log x$,
$\#\{p\le x:\ q_p<K\}\sim\sum_{p\le x}\min(1,K/p)$---the fixed-$K$
case down to $K=1$ is what governs the zero event and does NOT follow
from (i), since $\{q_p=0\}$ is a target of shrinking measure $1/p$;
\emph{(iii)} [Wieferich count, fluctuation-calibrated, with the
probability space declared]
$W(x)=\#\{\text{Wieferich } p\le x\}$ is a deterministic function,
so its fluctuation clauses are stated in the two legitimate forms,
an unqualified CLT for a fixed sequence having no meaning: the
\emph{deterministic envelope conjecture}
\[
\limsup_{x\to\infty}
\frac{W(x)-\log\log x}
{\sqrt{2\log\log x\,\log\log\log\log x}}\;=\;1
\]
(the iterated-logarithm scale at variance $V(x)\sim\log\log x$; a
bound $O_\varepsilon(V^{1/2+\varepsilon})$ would be weaker than
the model's exact prediction), and the \emph{almost-sure central
limit law, at the correct logarithmic weighting}: with
$u=\log\log x$ the variance-time of the count, sample $u$ with
density $du/u$ on $[\e,\,L]$, $L=\log\log X$ (equivalently: $\log u
=\log\log\log x$ uniform); then the sampled distribution of
$(W(x)-u)/\sqrt u$ converges to $N(0,1)$ as $X\to\infty$.  The
weighting is what makes the clause true, and the two natural
alternatives both fail.  A block-sampled version on $[X,X^A]$
fails because only $\log A=O(1)$ new events enter the block.  A
\emph{uniform}-in-$u$ version on $[1,L]$ fails too, by Brownian
scaling: with
$W(\e^{\e^u})-u\approx B(u)$ and $u=Ls$, the self-similarity
$B(Ls)/\sqrt{Ls}\overset{d}=\widetilde B(s)/\sqrt s$ means the
uniform-$u$ empirical distribution converges not to $\Phi$ but to
the \emph{random} occupation functional
$\int_0^1\mathbf 1\{\widetilde B(s)/\sqrt s\le z\}\,ds$---values at
proportional times stay correlated, so uniform sampling never
averages over enough independent scales.  The $du/u$ weighting is
the classical almost-sure-CLT weighting (the continuous analogue of
the $1/k$ weights), for which the model does predict almost-sure
convergence of the sampled law to $\Phi$.  Both clauses are
conjectured on the
Crandall--Dilcher--Pomerance side of the model choice that clause
(iv)'s discussion records.
\end{conjecture}

\setcounter{conjecture}{7}%
\begin{conjecture}[Uniform quantitative de Polignac; the core is
previously stated]\label{c16}
\emph{(i)} $\pi_d(x)=\#\{p\le x-d:\ p,\ p+d \text{ prime}\}
=\SG(d)\,\Li_2(x)\,(1+o(1))$ uniformly for even $d\le(\log x)^{A}$.
\emph{(ii)} [residual field, with derived covariance kernel]  The
same primes participate in many pair counts, and the covariance of
two of them is carried by triple configurations sharing one prime.
For even $d\neq d'$ define the overlap constant
\[
K(d,d')=\CT(0,d,d')+\CT(0,d',d+d')+\CT(0,d,d+d')
+\CT\bigl(0,\,|d-d'|,\,\max(d,d')\bigr),
\]
where $\CT$ denotes the Hardy--Littlewood triple constant of the
offset set (degenerate or inadmissible triples contributing $0$).
The field is randomized canonically---covariances of the
deterministic cumulative counts are otherwise undefined: for $t$ uniform on
$[x,2x]$ and a window length $H$, let
$\pi_d(t;H)=\#\{t<n\le t+H:\ n,\,n+d \text{ prime}\}$ and let
$Z_d(t;H)$ be its studentization.  The mean count per window is
$\asymp H/\log^2x$, and the regime is set by that ratio (a phase
condition: a bare $H\to\infty$ does not fix the regime):
\[
\operatorname{Cov}_t\bigl(\pi_d,\pi_{d'}\bigr)
\sim\frac{K(d,d')\,H}{\log^3x},
\qquad
\rho(d,d')\sim\frac{K(d,d')}{\sqrt{\SG(d)\SG(d')}\,\log x},
\]
valid for $\max(d,d')=o(H)$, which is the standing restriction of
this clause and holds throughout the window class fixed below.  In
general each overlap orientation $o$---the in-window displacements
$h_o\in\{0,\,d,\,-d',\,d-d'\}$ and their reflections---enters with
multiplicity $(H-|h_o|)_+$ rather than $H$, so that the shared-prime
term reads
\[
\frac1{\log^3x}\sum_{o}\bigl(H-|h_o|\bigr)_+\,K_3^{(o)}(d,d'),
\qquad \sum_o K_3^{(o)}(d,d')=K(d,d'),
\]
and an orientation with $|h_o|>H$ is absent from the window
altogether; under $\max(d,d')=o(H)$ one has $(H-|h_o|)_+=H(1+o(1))$
for every orientation and the displayed form is recovered, so the
uniform-in-$d$ statement carries that restriction throughout.
Here the shared-prime (triple) term is the \emph{first} covariance
contribution but not the whole of it: the full expansion is
\[
\operatorname{Cov}_t(\pi_d,\pi_{d'})
=\frac{K(d,d')\,H}{\log^3x}
+\frac1{\log^4x}\sum_{|h|\le H}(H-|h|)\,K_4(d,d';h)
+\text{error},
\]
with $K_4$ the centred singular series of the off-index four-point
sets $\{0,d,h,h+d'\}$---the cross analogue of the pinned diagonal
average of Conjecture~\ref{c12}---where degenerate offset relations
are handled inside the kernels: whenever the four-point set
collapses (repeated
entries, as at $h=0,\ h=d,\ h=-d'$, or $d'-d=\pm h$) the
configuration is a triple or pair and its contribution belongs to
the $K$-term (or to the mean), so $K_4$ is defined as the centred
singular series of \emph{genuinely four-element} sets, with special
pairs such as $d'=2d$ acquiring their extra overlap orientations in
$K(d,d')$'s census rather than double-counted across kernels.  The
triple term dominates only
where the pinned sum satisfies $\sum(1-|h|/H)K_4=o(\log x)$; the
dominance range is stated as an explicit function class: for
$H=\lambda(\log x)^{\alpha}$, any
fixed $\lambda>0$ and any $\alpha\ge1$ with $\alpha>A$---the
inequality on the exponents that makes $\max(d,d')=o(H)$ hold
uniformly for all $d,d'\le(\log x)^{A}$---the per-shift centred average
grows only polylogarithmically in $H$ (the Conjecture-\ref{c12}
scale), so $\sum(1-|h|/H)K_4/H=O((\log\log x)^{2+\varepsilon})
=o(\log x)$ and the triple term dominates throughout the
polylogarithmic window class; the conjecture asserts the two-term
form on exactly that class and takes no position outside it.  When
$H/\log^2x\to\lambda$ the
window counts converge jointly to Poisson laws with these vanishing
covariances, and when $H/\log^2x\to\infty$ every fixed finite
collection $\{Z_{d_1},\dots,Z_{d_r}\}$ is asymptotically jointly
normal with \emph{independent} limiting coordinates: both kernels
are finite-$x$ corrections to independence (test below), not
nondegenerate limiting kernels.  The diagonal correction to
$\operatorname{Var}_t(\pi_d)$ is Conjecture~\ref{c12}'s pinned
average, stated there.
\end{conjecture}

\setcounter{conjecture}{8}%
\begin{conjecture}[The sub-diffusivity of balanced prime races;
apparently new]\label{c18}
Model a balanced race (Conjecture~\ref{c21}(ii)'s symmetric class
differences) as the walk of its class-assignment steps, the process
defined at event index: enumerate
the pattern occurrences $n_1<n_2<\cdots$ in the two competing
classes, set $\xi_i=\pm1$ by class membership of $n_i$, and let
$\rho_k(x)$ denote the lag-$k$ sample autocorrelation of
$(\xi_i)$ over the events with $n_i\le x$.  Then the
walk is measurably \emph{sub-diffusive}:
\emph{(i)} [decaying-repulsion law] the step autocorrelations are
negative at small lags---the race analogue of the Lemke
Oliver--Soundararajan consecutive-pattern repulsion \cite{LOS},
which supplies the mechanism---and, since that repulsion is
scale-dependent rather than fixed, the conjectured form is a decay
law
\[
\rho_k(x)\;=\;-\,\frac{c_k\log\log x+d_k}{\log x}\,(1+o(1)),
\]
with constants $c_k,d_k$ in principle derivable from the LOS
correlation series (underived here; and the measured values
$\rho_1\approx-0.037$, $\rho_2\approx-0.009$ at $10^9$---universal
across twin and cousin races mod $5$ and $8$ to within $0.004$, each
${\sim}50$ standard errors from zero---are observations at one height,
consistent with $c_1\approx0.25$ there, not asymptotic constants);
\emph{(ii)} [long memory, measured] short lags do not exhaust the
deficit: the running maxima
$M(x)=\max_{t\le x}|D(t)|/\sqrt{\text{count}}$ of the four races
land at quantiles $0.02$--$0.17$ of the simulated independent-step
null (median $2.34$ at this span), and still only $0.03$--$0.23$
after correcting for $\rho_{1,2}$: the negative correlation persists
across lags, so the cumulative diffusivity
$\sigma^2(x)=1+2\sum_k\rho_k(x)$ sits well below $1$ at this
height.  The model layer is declared:
$(\xi_i)$ is a deterministic, nonstationary family and the
$\rho_k(x)$ are empirical statistics of it, so the Green--Kubo
identity $\sigma^2=1+2\sum_k\rho_k$ is invoked under an explicit
\emph{local-stationarity model hypothesis}---the step process
restricted to a window $[x,2x]$ is modelled as a stationary sequence
whose autocorrelations are the measured $\rho_k(x)$---and the
variance statement can equivalently be taken as a direct
\emph{definition} of $\sigma^2(x)$ through block averages of
$(S_{n+m}-S_n)^2/m$ over the events up to $x$, a form in which no
stationarity is assumed at all.  The Green--Kubo sum requires more than the
fixed-lag law
of (i): the
per-lag decay controls each $\rho_k(x)$ but not the sum unless the
decay is \emph{uniform in $k$ with summable tail}, since the number
of relevant lags may grow with $x$.  We therefore state the needed
hypothesis as part of the clause: $\rho_k(x)$ obeys the law of (i)
uniformly for $k\le K(x)$ with
$\sum_{k>K(x)}|\rho_k(x)|=o(\log\log x/\log x)$ for some
$K(x)\to\infty$, \emph{and}---since the tail bound alone leaves the
growing head sum uncontrolled---the
repulsion constants are summable, $\sum_k c_k<\infty$ and
$\sum_k|d_k|<\infty$, so that
$\sigma^2(x)=1-2\bigl(\textstyle\sum_kc_k\,\log\log
x+\sum_kd_k\bigr)/\log x\,(1+o(1))$ is a genuine consequence
rather than an unjustified interchange of limits.  The measured
correlation length at $10^9$ (a few tens of lags, with
$|\rho_k|$ decreasing geometrically in the observable range) is
consistent with both hypotheses.  Under them,
$\sigma^2(x)$ returns to $1$
asymptotically, making sub-diffusivity a finite-height phenomenon
of $\log\log x/\log x$ size, like every other second-order law in
this paper;
\emph{(iii)} [registered dichotomy, open] asymptotically the two
possibilities are \emph{mutually exclusive}, and which of them holds
is the open question.  Either the fixed-lag expansion of (i) extends
uniformly in $k$ with summable constants---in which case the
hypotheses of (ii) hold, (ii) applies, and the race is asymptotically
diffusive with $\sigma^2(x)\to1$ (running maxima
$\sim\sqrt{2\log\log x}$, Darling--Erd\H os class)---or that
expansion fails at long lags, and only in that case can the race
inherit the almost-periodic rigidity that Rubinstein--Sarnak-type
structure \cite{RubinsteinSarnak} forces on classical races (maxima
$\asymp(\log\log\log x)^{A}$, the Montgomery--Ng class).  Rigidity
therefore requires a failure of the hypotheses of (ii) at long lags,
and cannot coexist with them.  No explicit
formula is known for twin patterns, the two branches differ only
beyond any computable height, and we register the dichotomy without
choosing---what is conjectured is the finite-height law (i)--(ii).
\end{conjecture}

\section{Family laws and second-order refinements}

\setcounter{conjecture}{9}%
\begin{conjecture}[Uniform quadratic de Polignac; family law apparently
unstated, the $d=2$ instance previously stated, cf.\ OEIS A080149,
\cite{Carella}]\label{c1}
For even $d\ge2$ let
$\pi^{\mathrm q}_d(N)=\#\{n\le N:\ n^2+1,\ n^2+1+d \text{ both
prime}\}$ and $C(d)=C(x^2+1,\,x^2+1+d)$.  Then:
\emph{(i)} [uniform family law] for every fixed $B>0$,
$\pi^{\mathrm q}_d(N)=C(d)\,I_d(N)(1+o(1))$ uniformly over even
$d\le(\log N)^{B}$;
\emph{(ii)} [registered rate, strictly stronger] there is
$\eta_B>0$ with
\[
\sup_{\substack{2\le d\le(\log N)^{B}\\ 2\mid d}}
\Bigl|\frac{\pi^{\mathrm q}_d(N)}{C(d)\,I_d(N)}-1\Bigr|
\;\ll_B\;(\log N)^{-\eta_B}
\]
---an error clause that no version of Bateman--Horn supplies, registered
separately from (i) since its analytic content is independent;
\emph{(iii)} [family statistic: the law of the constants]
each local factor of $C(d)$ depends only on $d\bmod p$, and the
factors at any \emph{finite} set of primes are exactly independent
as $d$ varies (CRT), with the passage to the infinite product
controlled by the convergent tail-variance sum---which supplies the
uniform integrability (via $L^2$-boundedness of the normalized
partial products) that the moment identities require, since
almost-sure convergence alone would not carry means through the
limit; so the moments of
$C(d)$ over even $d\le D$ converge, as $D\to\infty$, to
\emph{derived} Euler products; in particular
\[
\frac1{\#}\sum_{2\mid d\le D}C(d)\to\bar C=2.7447\ldots,
\qquad
\operatorname{sd}\bigl(C(d)\bigr)\to1.6835\ldots,
\]
and the empirical distribution of $C(d)$ converges to the law of the
random product $\prod_p f_p(U_p)$ with independent uniform residues
$U_p$.  In particular ($d=2$, the classical instance):
$\#\{n\le N:\ n^2+1,\ n^2+3 \text{ both prime}\}\sim C(2)\,I(N)$,
$C(2)=2.954014\ldots$
\end{conjecture}

\setcounter{conjecture}{10}%
\begin{conjecture}[First occurrences of prime gaps, liminf form; the
family is classical, the quantified form apparently
unstated]\label{c19}
Let $\mathcal G$ be the set of even $g$ that occur as gaps between
consecutive primes (all even $g$, under Polignac's conjecture), and for
$g\in\mathcal G$ let $p(g)$ be the prime starting the first such gap.
Then:
\emph{(i)} $\log p(g)/\sqrt g$ is bounded between positive constants
on $\mathcal G$;
\emph{(ii)} [liminf law, stated as the dual of the maximal-gap
limsup, \emph{conditional on the realization hypothesis below}]
\[
\liminf_{\substack{g\to\infty\\ g\in\mathcal G}}
\ \frac{\log p(g)}{\sqrt g}
\;=\;\sqrt{\e^{\gamma}/2}\;=\;0.943682\ldots
\]
The inversion of Granville's $\limsup G(X)/\log^2X=2\e^{-\gamma}$
into first-occurrence coordinates is \emph{not} automatic: the limsup
controls how large maximal gaps get,
not which gap \emph{sizes} are realized as first occurrences near the
envelope.  Clause (ii) is therefore conditional on the
\emph{realization hypothesis}: infinitely many $g\in\mathcal G$ first
occur inside a maximal-gap event at the Cram\'er--Granville envelope
scale, i.e.\ with $g=(2\e^{-\gamma}+o(1))\log^2 p(g)$.  Under that
hypothesis the displayed value is forced by algebra; without it only
the lower bound
$\liminf\log p(g)/\sqrt g\ge\sqrt{\e^{\gamma}/2}$ follows from the
envelope.  The value is registered \emph{as} that dual (against the
slope-1 laws), not claimed as an independent phenomenon;
\emph{(iii)} [second-order, singular-series dependence] sampling $g$
uniformly from the realized gaps $\mathcal G\cap[G,2G]$ and letting
$G\to\infty$,
\[
\log p(g)\;=\;\sqrt g+\tfrac12\log g-\tfrac12\log\SG^{*}(g)+E_g,
\qquad
\SG^{*}(g)=\prod_{\substack{p\mid g\\ p>2}}\frac{p-1}{p-2},
\]
where the error $E_g$ converges in distribution: for the sampled $g$,
\[
\Pr\bigl[E_g\le t\bigr]\;\longrightarrow\;
1-\exp\bigl(-\e^{\,2(t-\mu)}\bigr)
\qquad(G\to\infty)
\]
for some centring constant $\mu$---the min-type (reversed) Gumbel
law \emph{at scale} $\tfrac12$, i.e.\ the law of $\tfrac12\log W$
for $W$ a unit exponential.  The scale is forced by the hazard computation:
the cumulative intensity of
gap-$g$ events is $\Lambda_g(y)\asymp(\SG(g)/y^2)\exp(y-g/y)$ in
$y=\log x$ (the factor $\e^{-g/y}$ \emph{is} the no-intervening-prime
probability), so
$\frac{d}{dy}\log\Lambda_g=1+g/y^2-2/y=2+o(1)$ at the centring
point $y_0\approx\sqrt g$, and the first-passage identity
$\Lambda_g(y)=W$ gives $E_g=y-y_0=\tfrac12\log W+o(1)$.  The scale is
$\tfrac12$ and not $1$ precisely because the slope of the
log-intensity at $y_0$ is $2$ (both the $y$-term and the $g/y$-term
contribute), and the same slope-$2$ expansion is what produces the
$\tfrac12$ coefficients of the centring.  (Gumbel-type modelling of gap
extremes and first occurrences is Kourbatov--Wolf territory
\cite{Kourbatov,KW2020}; the claim of this clause is only the
$\SG^{*}$-centring and the scale-$\tfrac12$ first-passage form,
not the extreme-value framework.  Sampling uniformly
over the \emph{realized}
gaps in
$[G,2G]$ conditions on realization, which the independent
first-passage model does not represent; the clause conjectures that
this selection does not shift the limiting law, and that assumption
is part of what a refutation would refute.)  Smooth
gaps, whose tuple constant is large, appear \emph{earlier} by exactly
half a logarithm of that constant---the waiting-time first-passage
computation makes the coefficient $-\tfrac12$, not a free parameter,
while $\mu$ is left as the one unresolved constant of the clause.  (The
restriction to $\mathcal G$ is needed: $p(g)$ is not known
to be defined for every even $g$, and quantifying over all $g$ would
smuggle Polignac's conjecture into the statement.  The
$\SG^{*}$-clause records that the Hardy--Littlewood factor of $g$
materially shifts waiting times; it does not move the liminf constant,
since $\log\SG^{*}(g)=O(\log\log g)$.)
\end{conjecture}

\setcounter{conjecture}{11}%
\begin{conjecture}[The Stern lane race; apparently new]\label{c14}
In Stern's representation problem $n=p+2k^2$ ($n$ odd), race the
$k$-even and $k$-odd lanes.  Prime-square contamination of the
$\Lambda$-weighted lane counts comes from $n=q^2+2k^2$---the norm
form $x^2+2y^2$ of $\mathbb Q(\sqrt{-2})$ with prime $x$---and since
$q^2\equiv1\pmod8$ while $2k^2\equiv0$ or $2\pmod 8$ by the parity of
$k$:
\emph{(i)} [class assignment, direct congruence] contamination
sits in the $k$-\emph{even} lane iff $n\equiv1\pmod8$, in the
$k$-\emph{odd} lane iff $n\equiv3\pmod8$, and for
$n\equiv5,7\pmod8$ \emph{no} square contamination exists at all:
$q^2+2k^2\equiv1$ or $3\pmod 8$ always, a two-line congruence (genus theory
enters only for the converse question of
which $n$ \emph{are} represented), so the null classes
are provable, not conjectured;
\emph{(ii)} [drift law, in the ensemble form of
Conjecture~\ref{c25}]  Fix the
ensemble: for a scale $X$, let
$\mathbb E_X$ average over odd $n\in(X,2X]$ of the given residue
class mod $8$ with weight proportional to $1/n$ (the logarithmic
sampling measure, as at Conjectures~\ref{c15} and~\ref{c25}).
Define, per class, the \emph{clean-minus-contaminated} difference
with explicit sign: $D(n)=R_{\mathrm o}-R_{\mathrm e}$ for
$n\equiv1\pmod 8$ (even lane contaminated),
$D(n)=R_{\mathrm e}-R_{\mathrm o}$ for $n\equiv3\pmod 8$, and
$D(n)=R_{\mathrm e}-R_{\mathrm o}$ (sign immaterial) on the null
classes $n\equiv5,7\pmod8$.  Set
$D_{\mathrm{sys}}(n)=\bigl[\sum_{(q,k):\,q^2+2k^2=n,\ q\
\mathrm{prime},\,k\ge1}\log q\bigr]/\lambda(n)$---the weight is
$\Lambda(q^2)=\log q$, each representation counted once in the
$k\ge1$ convention (the ordered-representation factor $2$ of
Conjecture~\ref{c25} has no counterpart here, and doubling would be a
normalization error)---where
$\lambda(n)$ is the \emph{analytic} lane mean of $\log p$ over the
representation profile ($\lambda(n)=\int\log(n-2u^2)\,d\mu_n$ under
the lane measure; the verification uses the empirical lane mean as its
estimator).  The law, coefficient-identifying as at
Conjecture~\ref{c25}(i): on the contaminated classes
\[
\mathbb E_X\bigl[D-D_{\mathrm{sys}}\bigr]
=o\bigl(\mathbb E_X[D_{\mathrm{sys}}]\bigr)
\qquad(X\to\infty),
\]
and on the null classes---where $D_{\mathrm{sys}}$ vanishes
identically, so that a comparison against it would be empty---the
statement is made against the positive comparison scale
\[
A(X)\;:=\;\mathbb E_X^{(1)}\bigl[D_{\mathrm{sys}}\bigr],
\]
the logarithmic-ensemble mean of the systematic census taken over the
\emph{contaminated} class $n\equiv1\pmod 8$: writing
$D_{\mathrm{null}}$ for the difference $D$ on a null class
$n\equiv5$ or $7\pmod 8$, the claim is
$\mathbb E_X[D_{\mathrm{null}}]=o\bigl(A(X)\bigr)$.
\end{conjecture}

\setcounter{conjecture}{12}%
\begin{conjecture}[Microscopic variance law; classical family, the
constant is Montgomery--Soundararajan's \cite{MS}]\label{c20}
Fix $\lambda>0$ and $c\in(0,1]$, and let $X$ count primes in
$[t,\,t+\lambda\log x)$ for $t$ uniform on $[x,(1+c)x]$.  Then, under
the Hardy--Littlewood $k$-tuple conjectures,
$X\Rightarrow\mathrm{Poisson}(\lambda)$ (Gallagher's argument
\cite{Gallagher}, which is conditional on exactly those conjectures),
and---\emph{conditional on the quantitative hypothesis} spelled out
after the statement (uniform Hardy--Littlewood with $o(h)$ control
of the averaged singular series), without which the display
below is a registered extrapolation rather than a consequence---at
finite $x$
\[
\frac{\operatorname{Var}X}{\mathbb E\,X}
\;=\;1-\frac{\log(\lambda\log x)+\gamma+\log2\pi-1}{\log x}
+o\!\Bigl(\frac1{\log x}\Bigr),
\]
with $\gamma+\log2\pi-1=1.41509\ldots$, the leading correction being
independent of the sampling constant $c$.  [Interpolation clause]
More generally, for windows of length $H=\lambda(\log x)^{\alpha}$
the \emph{same} formula with $\log H$ in place of
$\log(\lambda\log x)$,
\[
\frac{\operatorname{Var}X}{\mathbb E\,X}
\;=\;1-\frac{\log H+\gamma+\log2\pi-1}{\log x}
+o\!\Bigl(\frac1{\log x}\Bigr),
\]
holds uniformly for $1\le\alpha\le A$ (any fixed $A$ with
$H\le x^{o(1)}$): one expression interpolates from Gallagher's
Poisson boundary ($\alpha=1$, deficit $\asymp\log\log x/\log x$)
through the mesoscopic Montgomery--Soundararajan regime (deficit
$\asymp\alpha\log\log x/\log x$), with no transition other than the
slow growth of $\log H$.
\end{conjecture}

\setcounter{conjecture}{13}%
\begin{conjecture}[The null-mechanism race; apparently new]\label{c8}
For the quadratic twin pairs of Conjecture~\ref{c1} at $d=2$, square
contamination is \emph{algebraically impossible}: $n^2+1=q^2$ has no
solution with $n\ge1$, and $n^2+3=q^2$ only at $(n,q)=(1,2)$.  The
surviving classes of $n$ mod $5$ are $\{0,1,4\}$, and the classes
$1$ and $4$ have identical singular series (elementary, via CRT).
Let $\pi^{\mathrm q}_a(x)$ count solutions with $n\equiv a\pmod5$,
$D(x)=\pi^{\mathrm q}_1(x)-\pi^{\mathrm q}_4(x)$, and define the
normalized race process
$Y(x)=D(x)\big/\sqrt{\pi^{\mathrm q}_1(x)+\pi^{\mathrm q}_4(x)}$.
Then, in contrast to the twin race of Conjecture~\ref{c21}, this race
is \emph{driftless}:
\emph{(i)} $\M_x\bigl(D(t)\log^2t/\sqrt t\bigr)\to0$ at the
\emph{drift-scale} normalization of Conjecture~\ref{c21}(i)---the
scale at which the contaminated races converge to nonzero
constants---so this is the genuine negation of
Conjecture~\ref{c21}(i) for this pair (the weaker normalization
$\M_x(Y)\to0$ with $Y=D/\sqrt{\pi^{\mathrm q}}$ would not exclude a
drift at the contamination scale, hence would not make this race a
control; the form stated excludes exactly that);
\emph{(ii)} [occupation law, stated as a conditional consequence] the
arithmetic content of this clause is a \emph{strong approximation
hypothesis}, for the process defined precisely at event index: let
$n_1<n_2<\cdots$ enumerate the $n$ with $n^2+1$ and $n^2+3$ both
prime and $n\equiv1$ or $4\pmod5$, set $\xi_i=+1$ if
$n_i\equiv1$ and $-1$ if $n_i\equiv4$, and let
$S_N=\xi_1+\cdots+\xi_N$ (so $D(x)=S_{N(x)}$ with $N(x)$ the number
of such events to $x$).  The hypothesis is an \emph{almost-sure
invariance principle} (ASIP): on a common probability space carrying
$(\xi_i)$ and a standard Brownian motion $B$,
\[
S_N\;=\;B(N)+O\bigl(N^{1/2-\delta}\bigr)
\qquad\text{almost surely, for some }\delta>0.
\]
Weak convergence alone would not suffice, since a distributional
limit on each fixed interval does not control pathwise occupation
across unboundedly many logarithmic scales; under the ASIP the
almost-sure error is absorbed at every logarithmic scale and the
Lamperti/ergodic occupation argument below is valid.  The weak
functional limit theorem---$S_{\lfloor N\tau\rfloor}/\sqrt N$
converging weakly to standard Brownian motion on
$\tau\in[0,1]$---is retained only as the weaker, separately
falsifiable consequence of the ASIP.  Everything after that is
probability, not number theory: granting the ASIP,
$D(t)=B(N(t))+O\bigl(N(t)^{1/2-\delta}\bigr)$
with $N(t)\asymp t/\log^2t$, the error being $o(\sqrt{N(t)})$
almost surely, and in \emph{logarithmic} time
$u=\log t$ the
Lamperti reduction $Z(s)=\e^{-s/2}B(\e^s)$ is a stationary, ergodic
Ornstein--Uhlenbeck process, so the logarithmic occupation measure of
leadership converges almost surely to $\tfrac12$:
\[
\frac1{\log x}\int_2^x\mathbf 1_{\{D(t)>0\}}\,\frac{dt}t
\;\longrightarrow\;\frac12,
\]
with fluctuations of the Gaussian scale
$\sqrt{\log2/\log x}$ (the occupation-indicator covariance
$\frac1{2\pi}\arcsin\e^{-u/2}$ integrates to $\tfrac12\log2$; note
that the constant for the \emph{sign} average $2\mathbf 1-1$ is twice
this, the occupation fraction itself carrying
$\sqrt{\log2/\log x}$); the classical
arcsine spread survives only in the \emph{natural} event clock,
i.e.\ for occupation fractions of $\{B(m)>0,\ m\le N(x)\}$.  A
persistent systematic leader---a log-occupation stuck near $0$ or
$1$ beyond the stated Gaussian scale---is the refutation.
\end{conjecture}

\setcounter{conjecture}{14}%
\begin{conjecture}[The cubic shift family: the distribution of its
constants; instance-level law apparently unstated]\label{c2}
For non-cube $a\ge1$ let $C(a)=C(x^3+a)$.  Only $p\equiv1\pmod3$
moves the product, and the local root count there takes \emph{three}
values as $a$ varies mod $p$:
$\omega=3$ with probability $(p-1)/3p$ ($-a$ a nonzero cube),
$\omega=1$ with probability $1/p$ (the case $p\mid a$, easy to
overlook and essential to the mean-value identity below), and
$\omega=0$ with probability
$2(p-1)/3p$; for $p\equiv2\pmod 3$ the cube map is a bijection and
$\omega=1$ identically.  Then:
\emph{(i)} [mean-value one, elementary lemma]
$\mathbb E_a[\omega(p)]=3\cdot\frac{p-1}{3p}+\frac1p=1$ exactly at
every $p$, so the derived mean of $C(a)$ is exactly $1$---the cubic
analogue of the Korevaar--te Riele mean-value-one theorem for linear
prime-pair constants (arXiv:0806.1667), here a one-line computation
at each finite level; passing the mean through the infinite product
needs more than almost-sure convergence, and the needed ingredient is
on hand: the convergent
variance sum makes the partial products an $L^2$-bounded martingale
(each factor has mean $1$ and is independent of the earlier ones),
so they are uniformly integrable and the limit law has mean exactly
$1$, not merely limit-of-means $1$;
\emph{(ii)} [limit law] as $a\to\infty$ over non-cubes, $C(a)$
converges in distribution to the random Euler product
$\prod_{p\equiv1(3)}f_p(\xi_p)$ with the three-state local variables
$\xi_p$ above---independent at any finite set of primes exactly, by
CRT, with the infinite product's tail controlled by the convergent
variance sum (unlike the quadratic family of
Conjecture~\ref{c1}(iii), no normalization is needed)---with derived
standard deviation $0.2762\ldots$ computed from the full three-state
law;
\emph{(iii)} [uniformity] $\#\{n\le N: n^3+a \text{ prime}\}
=C(a)I_a(N)(1+o(1))$ uniformly over non-cube $a\le(\log N)^B$.
\end{conjecture}

\setcounter{conjecture}{15}%
\begin{conjecture}[Conjecture F as a family; the core is previously
stated, \cite{HL1923}, cf.~\cite{JW}]\label{c24}
For odd $A$ let $Q_A(N)=\#\{n\le N:\ n^2+n+A \text{ prime}\}$ and
$C(A)=C(x^2+x+A)$.  Then, for every fixed $B>0$,
$Q_A(N)=C(A)\,I_A(N)\,(1+o(1))$ uniformly over odd
$1\le A\le(\log N)^B$, with the residual field obeying the analogue
of Conjecture~\ref{c16}(ii), kernel derived and probability space
specified: for $t$ uniform on $[N,2N]$ and window length $H\le
N^{o(1)}$, let
$Q_A(t;H)=\#\{t<n\le t+H:\ n^2+n+A \text{ prime}\}$ (mean
${\asymp}H/2\log N$, which sets the regime).  Two members share the
index window, so with $C(A,A')=C(x^2+x+A,\ x^2+x+A')$,
\[
\operatorname{Cov}_t\bigl(Q_A,Q_{A'}\bigr)\;\sim\;
\bigl[C(A,A')-C(A)\,C(A')\bigr]\,\frac{H}{4\log^2N},
\]
with $C(A,A')=0$ (hence \emph{negative} correlation) for the locally
exclusive pairs.  This is the \emph{same-index} contribution only: the full
covariance adds the off-index pinned sum
\[
\frac1{4\log^2N}\sum_{0<|h|\le H}(H-|h|)\,
\bigl[C(f_A(x),f_{A'}(x+h))-C(A)C(A')\bigr],
\]
the two-parameter analogue of
Conjecture~\ref{c12}'s $G(H)$, whose evaluation is this family's
open component.  (The order of this sum matters: the joint event here is
\emph{two single prime values} at
distinct indices, of probability $\asymp1/4\log^2N$, unlike
Conjecture~\ref{c16}'s pair-of-pairs events at $1/\log^4x$.  At that
order the off-index sum, which grows like a pinned
$G$-average over $\asymp H$ shifts against the same-index term's
single $H$, is generically the \emph{larger} component, so the
numerical cancellation of the same-index kernel establishes nothing
about the full cross-member covariance, and the localization of the
observed profile deficit to the diagonal is an open question rather
than an inference.)  For $H/\log N\to\lambda$ the counts are jointly
Poisson with local exclusions, and for $H/\log N\to\infty$ each
fixed finite collection of studentized counts is asymptotically
jointly normal, all correlations of size $\asymp1/\log N$ making
both kernels finite-$N$ corrections to independence, exactly as at
Conjecture~\ref{c16}(ii).
(Uniformity over a range growing with $N$ is the substantive claim;
uniformity over any \emph{fixed} finite set of $A$ is logically
automatic from the individual asymptotics.)  Among
odd $A\le199$, Euler's $A=41$ has the largest constant,
$C(41)=6.6395463\ldots$ (Cohen's high-precision value \cite{Cohen}).
\end{conjecture}

\setcounter{conjecture}{16}%
\begin{conjecture}[Twin-member Goldbach, orientation-resolved; the
basis conjecture and the integral kernel are Dubner's
\cite{Dubner2000}; the orientation decomposition apparently
new]\label{c17}
For even $n$ let $R_T(n)$ count ordered $(a,b)$, $a+b=n$, with $a$
and $b$ both members of twin pairs, each ordered role (lower/upper)
counted.  Each of the four role assignments is a $4$-form linear
system in $t=a$ with root set
$\{0,\pm2\}\cup\{n,n\pm2\}$-type mod each $p$, hence its own
$n$-dependent singular series $\SG_4^{(o)}(n)$, and
\[
R_T(n)\;=\;\Bigl[\sum_{o}\SG_4^{(o)}(n)\Bigr]
\int_5^{n-5}\frac{dt}{\log^2t\,\log^2(n-t)}\;\bigl(1+o(1)\bigr),
\]
the four constants computed exactly per $n$ from the coincidence
pattern of the root sets (which of $n,n\pm2,n\pm4$ each small prime
divides).  The error term is stated as $(1+o(1))$, not
$(1+O(1/\log n))$: our own measurements show a second-order deficit
whose $1/\log n$ coefficient is not yet pinned down (below).  Two precision
notes: the
asymptotic runs over \emph{varying} $n$, so what is assumed is a
Hardy--Littlewood estimate for $4$-form systems \emph{uniform in the
$n$-dependent coefficients}---strictly stronger than the fixed-system
conjecture---and the number $5$, the unique prime that is
simultaneously a lower twin member (of $(5,7)$) and an upper one (of
$(3,5)$), is double-counted consistently by the ordered-role
convention, affecting finitely many representations per $n$ and no
asymptotic.
\end{conjecture}

\section{Instances and structural companions}

\setcounter{conjecture}{17}%
\begin{conjecture}[Contamination in prime triplets; apparently
new]\label{c3}
The contamination calculus of Conjecture~\ref{c21}(v) extends to
$k$-tuples, and the triplet pattern $(n,\,n+2,\,n+6)$ is its first
instance beyond pairs.  Of the three configurations placing a prime
square inside the pattern, two die algebraically for
$q>3$---$(q^2,q^2+2,q^2+6)$
by $3\mid q^2+2$, and $(q^2-6,q^2-4,q^2)$ because
$q^2-4=(q-2)(q+2)$ is composite (at $q=3$ the factorization is
trivial, $q^2-4=5$ is prime, and the single bounded exceptional
configuration $(3,5,9)$ occurs---one term, harmless to every
asymptotic)---leaving the single
\emph{doubly-thinned} survivor $(q^2-2,\ q^2,\ q^2+4)$, which
requires $q^2-2$ \emph{and} $q^2+4$ simultaneously prime and forces
$q\equiv\pm2\pmod5$ for $q\neq5$ (else $5\mid q^2+4$; at $q=5$ all
three conditions do hold, the configuration being $(23,25,29)$, but
its start $23\equiv3\pmod5$ lies outside the triplet start classes
$\{1,2\}$, so the single term is harmless).  Triplet starts lie in
classes $n\equiv1,2\pmod 5$, and the surviving configuration has
start $q^2-2\equiv2$: class $2$ is contaminated and class $1$ leads,
in the drift-scale form of Conjecture~\ref{c21}(i) (a normalization
at the noise scale, $\M_x((D-T_3)/\sqrt{\pi_3})$, would be vacuous
for the coefficient, exactly as at the pair races):
\[
\M_x\!\left(
\frac{\bigl(\pi_3(t;5,1)-\pi_3(t;5,2)\bigr)\log^3t}{\sqrt t}\right)
\to c_3:=\lim_{x\to\infty}\frac{T_3(x)\log^3x}{\sqrt x},
\quad
T_3(x)=\frac1{\log^3x}\!\!
\sum_{\substack{q\le\sqrt x,\ q\neq5\\ q^2-2,\ q^2+4\ \mathrm{prime}}}
\!\!\log(q^2-2)\log q\,\log(q^2+4)
\]
($c_3$ exists and is the Bateman--Horn constant of the triple
$(q,\,q^2-2,\,q^2+4)$, since the census sum is $\sim c_3\sqrt x$;
convergence of the logarithmic mean at this normalization is the
same Rubinstein--Sarnak-type structure hypothesis as at
Conjecture~\ref{c21}(i)), with drift-to-noise
$\asymp\log^{-3/2}x$---one logarithm weaker than
the pair races, a scale the calculus itself predicts and the
verification must respect.
\end{conjecture}

\setcounter{conjecture}{18}%
\begin{conjecture}[The contamination matrix: sexy pairs with two
surviving orientations; apparently new]\label{c7}
For the pattern $(n,n+6)$, \emph{both} prime-square orientations
survive, on complementary classes of $q$ mod $5$---the first matrix
instance of the calculus (every pair race so far had exactly one):
orientation $A=(q^2-6,\,q^2)$ requires $q^2-6$ prime, forcing
$q\equiv\pm2\pmod5$ for $q\neq5$ (at $q=5$ the value $19$ is prime,
but its start $19\equiv4\pmod5$ lies outside the sexy start classes
$\{1,2,3\}$), and lands in start class $3$ (mod $5$), class
$3$ (mod $8$); orientation $B=(q^2,\,q^2+6)$ requires $q^2+6$ prime,
forcing $q\equiv\pm1\pmod5$ (the term $q=5$ falls outside the start
classes), and lands in class $1$ (mod $5$), class $1$ (mod $8$).
Sexy starts occupy classes $\{1,2,3\}$ mod $5$ and $\{1,3,5,7\}$ mod
$8$; with
$T_{A}(x)=\frac1{\log^2x}\sum_{q^2-6\ \mathrm{prime},\,q\neq5}\log
q\log(q^2-6)$,
$T_{B}(x)=\frac1{\log^2x}\sum_{q^2+6\ \mathrm{prime},\,q\neq5}\log
q\log(q^2+6)$, and $c_A,c_B$ their drift-scale limits
($c_A=\lim T_A\log^2x/\sqrt x$, the Bateman--Horn constant of
$(q,q^2-6)$, likewise $c_B$ for $(q,q^2+6)$), the predicted drift
vector, in the drift-scale limiting-log-mean sense of
Conjecture~\ref{c21}(i) (an $\M_x$-over-$\sqrt\pi$ normalization would
be coefficient-blind), is:
mod 5: $\M_x\bigl((\pi(2)-\pi(3))\log^2t/\sqrt t\bigr)\to c_A$,
$\M_x\bigl((\pi(2)-\pi(1))\log^2t/\sqrt t\bigr)\to c_B$ (class $2$
\emph{clean}); mod 8:
$\M_x\bigl((\tfrac12(\pi(5)+\pi(7))-\pi(3))\log^2t/\sqrt t\bigr)\to
c_A$,
$\M_x\bigl((\tfrac12(\pi(5)+\pi(7))-\pi(1))\log^2t/\sqrt t\bigr)\to
c_B$ (classes $5,7$ clean, mutually symmetric, with
$\M_x\bigl((\pi(5)-\pi(7))\log^2t/\sqrt t\bigr)\to0$).  Two
independently computable drift constants,
one certified null class per modulus.
\end{conjecture}

\setcounter{conjecture}{19}%
\begin{conjecture}[Fibonacci--Lucas twins: the convergent side
stress-tested; the object is OEIS A080327]\label{c9}
\emph{(i)} [rank-disjointness, elementary, with its exception stated]
every \emph{odd} prime factor $r$ of $L_p$ has rank of apparition
exactly $2p$ (from $z(r)\mid2p$, $z(r)\nmid p$, and $z(r)=2$ being
impossible), and every prime factor of $F_p$ ($p\neq5$ prime) has
rank exactly $p$: the odd divisor pools of the two numbers are
disjoint.  The prime $2$ is the unique exception: $2$ divides both
$F_p$ and $L_p$ precisely when $p=3$
($F_3=2$, $L_3=4$; $\gcd(F_p,L_p)=2$ iff $3\mid p$)---the
restriction to odd divisors is essential, the clause being false at
$p=3$ without it.  Disjoint pools remove shared-\emph{prime}
correlation but not all correlation: dependence can still enter
through the order structure of the recurrence, which is why (iii)
matters;
\emph{(ii)} [finiteness] only finitely many primes $p$ have $F_p$ and
$L_p$ simultaneously prime;
\emph{(iii)} [calibration clause, quantitative] the
naive joint accounting---probability
$(\e^{\gamma}\log p/(p\log\phi))^2$ per index---is contradicted at
the $3\times10^{-3}$ level by the record: it prices the catalogued
index
$p=148091$ (OEIS A080327; \emph{both} $F_p$ and $L_p$ are
\emph{probable} primes---numbers of roughly $30{,}950$ digits, with
Lucas numbers proved prime only up to index $56003$---so the
contradiction is conditional on both probable-prime classifications,
for which no Baillie--PSW pseudoprime is known: no unconditional
refutation is claimed from uncertified numbers) at
that much expected mass beyond $10^4$.  A single surprising event is
not a strict refutation, and it does not by itself estimate the
\emph{size} of the correction: an underestimated constant, genuine
cross-dependence,
and a mis-calibrated finite-range tail are competing explanations
that one event cannot separate.  What it does establish
(conditionally as above) is that the naive model is wrong somewhere,
and it is fatal
to any Goldilocks completeness list; the conjectural content of (ii)
is that even the corrected accounting has convergent sum.
\end{conjecture}

\setcounter{conjecture}{20}%
\begin{conjecture}[Factorial twins; the uniqueness core is previously
stated, OEIS A088054; clauses (i) and (iii) are ours]\label{c10}
\emph{(i)} [window rigidity, elementary] for $n\ge4$ and
$2\le|a|\le n$ the number $n!+a$ is composite (any prime factor $p$
of $a$ satisfies $p\le n$, so $p\mid n!+a$, and $n!+a>n\ge p$); the
two small anomalies $2!-2=0$ and $3!-3=3$ are the only exceptions,
which is why the clause is restricted to $n\ge4$.  Hence among all offsets
of bounded size the only
candidates for primality near $n!$ are $a=\pm1$: the twin question
below is the \emph{unique} bounded-offset constellation question at
$n!$, and every admissible offset set is a subset of $\{-1,+1\}$.
\emph{(ii)} [uniqueness; previously stated] Only finitely many $n$ have $n!-1$ and
$n!+1$ simultaneously prime, and the complete list is $n=3$: i.e.\
$6=3!=3\#$ is the unique factorial lying between twin primes.
\emph{(iii)} [joint fluctuation model] the two single-sided counts
$F_{\pm}(N)=\#\{n\le N:\ n!\pm1 \text{ prime}\}$ are each
$\sim\e^{\gamma}\log N$ and, \emph{under the independent-indices
model stated below}, asymptotically independent, with
$F_{+}-F_{-}$ normalized by $\sqrt{2\e^{\gamma}\log N}$
asymptotically standard normal.  The probability space is declared:
the clause is about a random model in which the events
$E_n^{\pm}=\{n!\pm1\ \text{prime}\}$ carry the Caldwell--Gallot
marginal hazards $\e^{\gamma}/n$ with an \emph{unspecified}
dependence structure, and the Gaussian limit is asserted under three explicit hypotheses on that structure: the same-index bound $\sum_{n\le N}\bigl|\operatorname{Cov}(\mathbf 1_{E_n^+},\mathbf 1_{E_n^-})\bigr|=o(\log N)$, which pins the variance at $2\e^{\gamma}\log N\,(1+o(1))$ and which the joint law $\sim(\e^{\gamma}/n)^2$ derived below supplies with a convergent sum---without it the remaining hypotheses leave the second cumulant free, and a model with $\Pr(E_n^+\cap E_n^-)=\e^{\gamma}/2n$ halves the limiting variance; the cross-index covariance bound
$\sum_{m<n\le N}\bigl|\operatorname{Cov}(\mathbf 1_{E_m},
\mathbf 1_{E_n})\bigr|=o(\log N)$ over the four sign pairs (the
variance itself is $\asymp\log N$), \emph{and}---since pairwise
control alone cannot preclude higher-order dependence---the
higher-cumulant condition
$\kappa_r(F_+-F_-)=o\bigl((\log N)^{r/2}\bigr)$ for each fixed
$r\ge3$, the method-of-moments sufficient condition.  Neither is
imported from an independence model; whether they hold for the
arithmetic dependence across
factorial indices is the open
component of this clause.
\end{conjecture}

\begin{theorem}[Cube obstruction]\label{t:cube}
For $k\ge2$, the cube $k^3$ is representable as $p+j^3$ with $p$ prime,
$j\ge1$, if and only if $3k^2-3k+1$ is prime.  Consequently the set of
integers not representable as $p+k^3$ is infinite: the non-representable
\emph{cubes} alone have counting function $\sim x^{1/3}$.  (That the
full exceptional set is $\sim x^{1/3}$ follows only when combined with
Conjecture~\ref{c13}(i), and is conjectural.)
\end{theorem}

\begin{proof}
$k^3-j^3=(k-j)(k^2+kj+j^2)$; for $1\le j\le k-2$ both factors exceed
$1$, so the difference is composite.  The only candidate is $j=k-1$,
giving $3k^2-3k+1$.  Since $3k^2-3k+1$ is composite for a density-one
set of $k$ (all but $O(K/\log K)$ of $k\le K$, by an upper-bound sieve
of Brun or Selberg type applied to the quadratic's prime values), all
but a vanishing proportion of cubes are unrepresentable.
\end{proof}

\setcounter{conjecture}{21}%
\begin{conjecture}[The boundary trichotomy for polynomial ladders; the
divisibility principle is classical and its cubic case is
Cunningham's cuban observation (OEIS A002407); the family
classification is apparently new]\label{c13}
For $F\in\mathbb Z[x]$ with $\deg F\ge2$ and \emph{positive leading
coefficient}, and $m>j\ge1$,
$(m-j)\mid F(m)-F(j)$ (textbook), and the divided difference
$(F(m)-F(j))/(m-j)$---whose leading form is then positive on the cone
$m>j\ge1$---exceeds $1$ outside an effectively bounded region; so
$F(m)-F(j)$ prime forces $m-j=1$ apart from finitely many
exceptional pairs (for the family below the exceptional set is
provably \emph{empty}).  Both hypotheses are necessary: for
$\deg F=1$ the cofactor can be identically $1$ and the collapse
fails, and for negative leading coefficient the divided difference
tends to $-\infty$ on the cone, so the ``exceeds $1$'' clause would be
false as stated (one may equivalently keep general sign and read the
claim through $|F(m)-F(j)|$).  Every representation problem $F(m)=p+F(j)$ in
this range
thus collapses to the boundary polynomial $D_F(m)=F(m)-F(m-1)$.  The content
is the classification
of the boundary lanes.  For $F=x^3+cx$, $c\ge0$:
\emph{(i)} [trichotomy, elementary] $D_F=3m^2-3m+1+c$ and the lane is
\emph{dead-parity} for $c$ odd ($D_F$ always even),
\emph{dead-3-adic} for even $c\equiv2\pmod3$ ($3\mid D_F$ always),
and admissible exactly for $c\equiv0,4\pmod6$;
\emph{(ii)} [uniform boundary law] over the admissible lanes,
$\#\{m\le M: D_F(m) \text{ prime}\}\sim C(D_F)\,I(M)$ uniformly for
$c\le(\log M)^B$---a Bateman--Horn family whose members are indexed
by the ladder classification;
\emph{(iii)} [attributed core] the case $c=0$ is the cuban ladder:
Cunningham's classical observation that a prime difference of cubes
forces consecutive arguments, with the non-cube exceptional set of
$n=p+k^3$ finite (Hardy--Littlewood's $E_3(X)=O(1)$ \cite{HL1923}).
\end{conjecture}

\setcounter{conjecture}{22}%
\begin{conjecture}[The power-obstruction ladder: an elementary
structural proposition with a Bateman--Horn corollary; apparently new
as a family]\label{c4}
For $k\ge2$ and $m\ge2$, write $D_k(m)=m^k-(m-1)^k$, and call $m^k$
representable if $m^k=p+j^k$ for some prime $p$ and $j\ge1$.  Then:
\emph{(i)} [theorem] for \emph{composite} $k$, no $k$-th power is
representable: if $r$ is a proper divisor of $k$ with $r>1$ (such
$r$ exists precisely because $k$ is composite; the excluded value
$r=1$ would allow only the useless factor $m-j=1$ at $j=m-1$), then
for every
$1\le j<m$ the difference $m^k-j^k$ has the proper factor $m^r-j^r>1$
with cofactor $>1$;
\emph{(ii)} [elementary] for \emph{prime} $k$ (including $k=2$),
$m^k$ is representable if and only if $D_k(m)$ is prime; and $D_k$ is
irreducible for prime $k$ (a root $\alpha$ has $(\alpha/(\alpha-1))^k=1$ with
$\alpha/(\alpha-1)\neq1$, so $\zeta=\alpha/(\alpha-1)$ is a primitive
$k$-th root of unity and $\mathbb Q(\alpha)=\mathbb Q(\zeta)$ has degree
$k-1=\deg D_k$, forcing irreducibility; explicitly
$D_k(x)=(x-1)^{k-1}\,\Phi_k(x/(x-1))$.  The forcing $j=m-1$ needs only
$k\ge2$);
\emph{(iii)} [conjecture] for each prime $k\ge3$ the representable lane follows
Bateman--Horn ($k=2$ is a theorem, the prime number theorem on the
lane $2m-1$):
$\#\{m\le M:\ D_k(m) \text{ prime}\}\sim C_k\int_2^M dt/\log D_k(t)$.
\end{conjecture}

\setcounter{conjecture}{23}%
\begin{conjecture}[The alternating cyclotomic chain; apparently
new]\label{c6}
There are infinitely many primes $p$ such that, writing
$u=\Phi_3(p)=p^2+p+1$, the three integers
\[
p,\qquad u=p^2+p+1,\qquad
\Phi_6(u)=u^2-u+1=p^4+2p^3+2p^2+p+1
\]
are all prime, with
$\#\{p\le x\}\sim C_6^{\mathrm{ch}}\,I(x)$ for the Bateman--Horn
system $\{x,\ x^2+x+1,\ x^4+2x^3+2x^2+x+1\}$,
$C_6^{\mathrm{ch}}=3.6143\pm0.011$.  The first chains are
$(2,7,43)$ and $(3,13,157)$.
\end{conjecture}

\setcounter{conjecture}{24}%
\begin{conjecture}[Twin cyclotomic bases: the complete quadratic
family; no prior trace found]\label{c5}
For the three cyclotomic polynomials of degree $2$ (indices $k$ with
$\varphi(k)=2$, i.e.\ $k\in\{3,4,6\}$), consider
$\Phi_k(n)$ and $\Phi_k(n+1)$ simultaneously prime.  Then:
\emph{(i)} [parity obstruction, elementary] for $k=4$ one member of
the pair $(n^2+1,\,n^2+2n+2)$ is even for \emph{every} $n$, so the
only prime pair is $n=1$, giving $(2,5)$---the same obstruction that
forbids consecutive $n^2+1$ primes;
\emph{(ii)} [collapse, elementary] $\Phi_6(x)=\Phi_3(x-1)$
identically, so the $k=6$ pair at index $n$ \emph{is} the $k=3$ pair
at index $n-1$: the family contains exactly one nontrivial statement;
\emph{(iii)} [conjecture, the live instance] there are infinitely many
$n$ for which the length-3 repunits in bases $n$ and $n+1$ are
simultaneously prime:
$\#\{n\le N:\ \Phi_3(n)=n^2+n+1 \text{ and }
\Phi_3(n+1)=n^2+3n+3 \text{ both prime}\}\sim C_5\,I(N)$,
with $C_5=C(x^2+x+1,\,x^2+3x+3)$.
\end{conjecture}

\begin{thebibliography}{99}\small

\bibitem{HL1923} G.~H. Hardy and J.~E. Littlewood,
\emph{Some problems of `Partitio numerorum'; III: On the expression of
a number as a sum of primes}, Acta Math.\ 44 (1923), 1--70.
DOI: 10.1007/BF02403921.

\bibitem{BH} P.~T. Bateman and R.~A. Horn,
\emph{A heuristic asymptotic formula concerning the distribution of
prime numbers}, Math.\ Comp.\ 16 (1962), 363--367.
DOI: 10.1090/S0025-5718-1962-0148632-7.

\bibitem{HB} D.~R. Heath-Brown,
\emph{Primes represented by $x^3+2y^3$},
Acta Math.\ 186 (2001), 1--84.  DOI: 10.1007/BF02392715.

\bibitem{CaldwellGallot} C.~K. Caldwell and Y.~Gallot,
\emph{On the primality of $n!\pm1$ and
$2\times3\times5\times\cdots\times p\pm1$},
Math.\ Comp.\ 71 (2002), 441--448.
DOI: 10.1090/S0025-5718-01-01315-1.

\bibitem{Lillie} G.~Lillie,
\emph{About the primality of primorials},
arXiv:2110.04302 [math.NT] (2021).

\bibitem{RubinsteinSarnak} M.~Rubinstein and P.~Sarnak,
\emph{Chebyshev's bias},
Experiment.\ Math.\ 3 (1994), no.~3, 173--197.
DOI: 10.1080/10586458.1994.10504289.

\bibitem{GranthamGranville} J.~Grantham and A.~Granville,
\emph{Fibonacci primes, primes of the form $2^n-k$ and beyond},
J.~Number Theory 261 (2024), 190--219; arXiv:2307.07894.
DOI: 10.1016/j.jnt.2024.02.002.

\bibitem{Wagstaffheu} S.~S. Wagstaff, Jr.,
\emph{Divisors of Mersenne numbers},
Math.\ Comp.\ 40 (1983), 385--397.
DOI: 10.1090/S0025-5718-1983-0679454-X.

\bibitem{Dubner2000} H.~Dubner,
\emph{Twin prime conjectures},
J.~Recreational Math.\ 30, no.~3 (1999--2000) (article reproduced at
OEIS A007534).

\bibitem{Dubner2002} H.~Dubner,
\emph{Carmichael numbers of the form $(6m+1)(12m+1)(18m+1)$},
J.~Integer Seq.\ 5 (2002), Article 02.2.1.

\bibitem{Kowalski} E.~Kowalski,
\emph{Averages of Euler products, distribution of singular series and
the ubiquity of Poisson distribution},
Acta Arith.\ 148 (2011), no.~2, 153--187; arXiv:0805.4682.

\bibitem{OSHP} T.~Oliveira e Silva, S.~Herzog, and S.~Pardi,
\emph{Empirical verification of the even Goldbach conjecture and
computation of prime gaps up to $4\cdot10^{18}$},
Math.\ Comp.\ 83 (2014), no.~288, 2033--2060.
DOI: 10.1090/S0025-5718-2013-02787-1.

\bibitem{Martin} K.~Martin,
\emph{Refined Goldbach conjectures with primes in progressions},
Exp.\ Math.\ 31 (2022), 226--232; arXiv:1806.00946.
DOI: 10.1080/10586458.2019.1596849.

\bibitem{GL} D.~A. Goldston and A.~H. Ledoan,
\emph{Jumping champions and gaps between consecutive primes},
Int.\ J. Number Theory 7 (2011), 1413--1421; arXiv:0910.2960.
DOI: 10.1142/S179304211100471X.

\bibitem{Kourbatov} A.~Kourbatov,
\emph{Maximal gaps between prime $k$-tuples: a statistical approach},
J.~Integer Seq.\ 16 (2013), Article 13.5.2; arXiv:1301.2242.  See also
arXiv:1309.4053.

\bibitem{Shanks} D.~Shanks,
\emph{On maximal gaps between successive primes},
Math.\ Comp.\ 18 (1964), 646--651.
DOI: 10.1090/S0025-5718-1964-0167472-8.

\bibitem{Granville} A.~Granville,
\emph{Harald Cram\'er and the distribution of prime numbers},
Scand.\ Actuar.\ J.\ 1995:1, 12--28.
DOI: 10.1080/03461238.1995.10413946.

\bibitem{Gallagher} P.~X. Gallagher,
\emph{On the distribution of primes in short intervals},
Mathematika 23 (1976), 4--9.  DOI: 10.1112/S0025579300016442.

\bibitem{MS} H.~L. Montgomery and K.~Soundararajan,
\emph{Primes in short intervals},
Comm.\ Math.\ Phys.\ 252 (2004), 589--617.
DOI: 10.1007/s00220-004-1222-4.

\bibitem{HLS} T.~Haddad, S.-K. Leung, and C.~Sabuncu,
\emph{Visiting early at prime times}, arXiv:2408.11781 (2024).
DOI: 10.48550/arXiv.2408.11781.

\bibitem{Wagstaff} S.~S. Wagstaff, Jr.,
\emph{Greatest of the least primes in arithmetic progressions having a
given modulus}, Math.\ Comp.\ 33 (1979), 1073--1080.
DOI: 10.1090/S0025-5718-1979-0528061-7.

\bibitem{CDP} R.~Crandall, K.~Dilcher, and C.~Pomerance,
\emph{A search for Wieferich and Wilson primes},
Math.\ Comp.\ 66 (1997), 433--449.
DOI: 10.1090/S0025-5718-97-00791-6.

\bibitem{JW} M.~J. Jacobson, Jr.\ and H.~C. Williams,
\emph{New quadratic polynomials with high densities of prime values},
Math.\ Comp.\ 72 (2003), 499--519.
DOI: 10.1090/S0025-5718-02-01418-7.

\bibitem{LOS} R.~J. Lemke Oliver and K.~Soundararajan,
\emph{Unexpected biases in the distribution of consecutive primes},
Proc.\ Natl.\ Acad.\ Sci.\ USA 113 (2016), E4446--E4454;
arXiv:1603.03720.  DOI: 10.1073/pnas.1605366113.

\bibitem{KW2020} A.~Kourbatov and M.~Wolf,
\emph{On the first occurrences of gaps between primes in a residue
class}, J.~Integer Seq.\ 23 (2020), Article 20.9.3; arXiv:2002.02115.

\bibitem{SL} M.-A. Sanchis-Lozano,
\emph{A heuristic study of the distribution of primes in short and
not-so-short intervals}, arXiv:1804.07659 (2018).
DOI: 10.48550/arXiv.1804.07659.

\bibitem{Sahoo} S.~Sahoo,
\emph{On twin prime distribution and associated biases},
arXiv:2111.09053 (2021).  DOI: 10.48550/arXiv.2111.09053.

\bibitem{PairBias} W.~Puszkarz,
\emph{Statistical bias in the distribution of prime pairs and isolated
primes}, arXiv:1807.00406 (2018).  DOI: 10.48550/arXiv.1807.00406.

\bibitem{Leung} S.-K. Leung,
\emph{Moments of primes in progressions to a large modulus},
Forum Mathematicum, published online 26 August 2024;
arXiv:2402.07941.

\bibitem{Fiori} A.~Fiori,
\emph{The least prime in arithmetic an progression} [sic],
arXiv:2404.02329 (2024).  DOI: 10.48550/arXiv.2404.02329.

\bibitem{Cohen} H.~Cohen,
\emph{High precision computation of Hardy--Littlewood constants},
unpublished notes (PDF hosted at OEIS A221712).

\bibitem{Carella} N.~A. Carella,
\emph{Twin primes in quadratic arithmetic progressions},
arXiv:1710.07827 (unrefereed preprint).
DOI: 10.48550/arXiv.1710.07827.

\end{thebibliography}
\end{document}